% from ega3-7-fr.tex
\oldpage[III]{215}
\section*{TABLE DES MATIÈRES}

\begin{flushleft}
\textsc{Chapitre III.} --- \textbf{Étude cohomologique des faisceaux cohérents
(suite)} \dotfill 5

\medskip
\S~6. Foncteurs Tor locaux et globaux ; formule de Künneth \dotfill 5

6.1. Introduction \dotfill 5

6.2. Hypercohomologie des complexes de Modules sur un préschéma \dotfill 6

6.3. Hypertor de deux complexes de modules \dotfill 9

6.4. Foncteurs hypertor locaux de complexes de Modules quasi-cohérents ; cas
des schémas affines \dotfill 14

6.5. Foncteurs hypertor locaux de complexes de Modules quasi-cohérents : cas
général \dotfill 16

6.6. Foncteurs hypertor globaux de complexes de Modules quasi-cohérents et
suites spectrales de Künneth : cas de la base affine \dotfill 21

6.7. Foncteurs hypertor globaux de complexes de Modules quasi-cohérents et
suites spectrales de Künneth : cas général \dotfill 25

6.8. Les suites spectrales d'associativité des hypertor globaux \dotfill 32

6.9. Les suites spectrales de changement de base dans les hypertor globaux
\dotfill 34

6.10. Structure locale de certains foncteurs cohomologiques \dotfill 39

\medskip
\S~7. Étude du changement de base dans les foncteurs homologiques covariants
de Modules \dotfill 43

7.1. Foncteurs de $A$-modules \dotfill 43

7.2. Caractérisation du foncteur produit tensoriel \dotfill 44

7.3. Critères d'exactitude des foncteurs homologiques de modules \dotfill 48

7.4. Critères d'exactitude pour les foncteurs
$\mathrm H_i(P\otimes_A M)$ \dotfill 53

7.5. Cas des anneaux locaux noethériens \dotfill 58

7.6. Descente des propriétés d'exactitude. Théorème de semi-continuité et
critère d'exactitude de Grauert \dotfill 60

7.7. Application aux morphismes propres : I. La propriété d'échange
\dotfill 65

7.8. Application aux morphismes propres : II. Critères de platitude
cohomologique \dotfill 72

7.9. Application aux morphismes propres : III. Invariance de la
caractéristique d'Euler-Poincaré et du polynôme de Hilbert \dotfill 76

\begin{flushright}
\emph{Reçu le 15 décembre 1962.}
\end{flushright}
\end{flushleft}

% La page imprimée 216 est blanche.
\clearpage
\thispagestyle{empty}
\null
\clearpage

